\section{挖角输入的真实三阶渐近}\label{cut:section}

有限尺度搜索提示四角挖去模型并未改善三阶系数。要把这个观察变为结论，
必须控制尺度离开一阶平坦区时的损失。下面的投影估计同时处理尺度并列和移动分支。

\begin{theorem}[紧挖角参数上的统一展开]\label{cut:main}
令 $Q_0=[-1/2,1/2]^2$，从其四角各挖去一个边长为 $kt$ 的方块，所得集合记为
$E_{k,t}$；边界取舍不影响测度。对每个 $K<\infty$，存在 $t_K>0$，使
$0\le k\le K$、$0<t<t_K$ 时
\begin{equation}\label{cut:expansion}
 \frac{|H_{1,\mathcal Q_c}(1-t,E_{k,t})|}{|E_{k,t}|}
 =2t^2+B(k)t^3+O_K(t^4),
\end{equation}
其中余项在 $k\in[0,K]$ 上一致，且
\[
 B(k)=\begin{cases}
 2-4k^3,&0\le k\le1/2,\\
 2+4(k-1)k^2,&1/2\le k\le2/3,\\
 6-18k+22k^2-8k^3,&2/3\le k\le1,\\
 -2+12k-8k^2,&k\ge1.
 \end{cases}
\]
特别地 $B(k)\le2$，等号恰在 $k=0,1$ 取得。结论允许 $k=k(t)$ 在该紧区间内变化。
\end{theorem}

\begin{proof}
首先 $|E_{k,t}|=1-4k^2t^2$，且 $|E_{k,t}\triangle Q_0|=4k^2t^2$。
原点处边长为 $1$ 的方体是高水平见证，因为
$\sqrt{1-4k^2t^2}>1-t$ 对充分小的 $t$ 一致成立；以下就用这个中心在原点的方体作为共同框架的锚点。
由共同框架引理~\ref{gl:frame}（先按实际质量伸缩，再伸缩回来），
全部高水平见证的中心为 $O_K(t)$，边长为 $1+O_K(t)$。
同一论证也适用于阈值 $1-Ct$，其中 $C$ 固定。

对有界 $v_1,v_2\ge0$、$a\in\mathbb R$，令中心为 $tv$、边长为 $1+ta$，并定义
\[
 b_j=\max(0,-a,v_j-a/2),\qquad
 U_j=[k+v_j+a/2]_0^k+[k-v_j+a/2]_0^k,
 \quad [s]_0^k=\min(k,\max(0,s)).
\]
逐坐标交长给出精确式
\[
 |E_{k,t}\cap Q|=(1-tb_1)(1-tb_2)-t^2U_1U_2.
\]
第一项是未挖角正方形的交叠；四个角块的交长乘积之和恰为第二项。
因所有参数有界，归一化后一致展开为
\begin{equation}\label{cut:local-expansion}
 \frac{|E_{k,t}\cap Q|}{(1+ta)\sqrt{|E_{k,t}|}}
 =1-tD(v,a)+t^2F_k(v,a)+O_K(t^3),
\end{equation}
其中
\[
 D=a+b_1+b_2,\qquad F_k=aD+b_1b_2-U_1U_2+2k^2.
\]
在任意固定有界参数框架内，$F_k$ 关于 $(v,a)$ 有统一 Lipschitz 常数。

令 $v=r(x,1-x)$，$0\le x\le1$，并记 $m=\min(x,1-x)$。
直接分段比较 $b_j$，有
\begin{equation}\label{cut:projection-gap}
 D(v,a)\ge r+\tfrac12\operatorname{dist}(a,I_{r,x}),
 \qquad I_{r,x}=[-2rm,2rm],
\end{equation}
且 $a\in I_{r,x}$ 时 $D=r$。把 $a$ 投影为 $a^\circ\in I_{r,x}$，
记距离为 $d$。由式~\eqref{cut:local-expansion}及 Lipschitz 界，
\[
 1-tD+t^2F_k(v,a)
 \le1-tr+t^2F_k(v,a^\circ)-td/2+Lt^2d.
\]
对 $t<1/(2L)$，最后两项之和非正。
反向直接使用平坦区中的尺度，便得到最大归一化平均为
\[
 1-tr+t^2\max_{a\in I_{r,x}}F_k(v,a)+O_K(t^3).
\]
这里只需比较值至少为 $1-Ct$ 的尺度：平坦区已经提供这种下界，
而上面的共同框架把全部这些尺度放入同一有界参数区间。
所以没有遗漏远处尺度，也没有交换不一致的极限和上确界。

平坦区内 $b_j=v_j-a/2$，且
$U_j=k+(k-v_j+a/2)_+$，因此在 $r=1$ 时
\begin{align*}
 G_k(x)&:=\max_{|a|\le2m}\Bigl\{
 a^2/4+a/2+x(1-x)+2k^2\\
 &\hspace{2.5em}{}-[k+(k-x+a/2)_+][k+(k-1+x+a/2)_+]\Bigr\}.
\end{align*}
所有这些函数在紧参数框架上一致 Lipschitz，平坦区端点也随 $r$ Lipschitz 变化。
先由上式的粗界得到径向边界 $r=1+O_K(t)$，再代入得到一致边界夹逼
\begin{equation}\label{cut:radial}
 r_{k,t}(x)=1+tG_k(x)+O_K(t^2).
\end{equation}
具体地，把 $r=1+tG_k(x)\pm Ct^2$ 代回平均展开，取足够大的固定 $C$，
即可分别得到严格高于和严格低于 $1-t$ 的值。
输入关于各坐标偶对称且逐坐标递减；与中心区间取交叠、再对尺度取上确界，
仍保持逐坐标递减。因此光晕在每条第一象限射线上是从原点开始的区间，
其端点满足上述夹逼，严格边界不影响面积。

为显式积分 $G_k$，由对称性只需 $0\le x\le1/2$。置 $u=a/2$，
关于 $u\in[-x,x]$ 的目标在折点 $x-k$、$1-x-k$ 之间分别凸或仿射，
故只需检查 $u=\pm x$ 及位于区间中的两个折点。
这些有限候选直接比较后得到下表；相邻区间端点的值相同。
\[
\begin{array}{c|c|c}
 k\text{ 的范围}&x\text{ 的范围}&G_k(x)\\\hline
 {}[0,1/2]&[0,(1-k)/2]&2x\\
 {}[0,1/2]&[(1-k)/2,1/2]&2x+2k(1-k-2x)\\\hline
 {}[1/2,2/3]&[0,(1-k)/2]&2x\\
 {}[1/2,2/3]&[(1-k)/2,1/2]&2[(1-k)^2+(2k-1)x]\\\hline
 {}[2/3,1]&[0,(1-k)/2]&2x\\
 {}[2/3,1]&[(1-k)/2,1-k]&2[(1-k)^2+(2k-1)x]\\
 {}[2/3,1]&[1-k,k/2]&2kx\\
 {}[2/3,1]&[k/2,1/2]&k(2k-1)+2(1-k)x\\\hline
 {}[1,\infty)&[0,1/2]&2k(1-k)+2(2k-1)x
\end{array}
\]
坐标变换 $(x,r)\mapsto tr(x,1-x)$ 的雅可比绝对值为 $t^2r$。
由式~\eqref{cut:radial}，全平面光晕面积为
\[
 2t^2\int_0^1r_{k,t}(x)^2\,dx
 =2t^2+4t^3\int_0^1G_k(x)\,dx+O_K(t^4).
\]
再除以 $1-4k^2t^2$ 只影响四阶及更高阶。逐行积分上述表格，
$B(k)=8\int_0^{1/2}G_k(x)\,dx$ 正是定理中的分段式。
前两段显然不超过 $2$；第三段的差为
$2-B=2(k-1)(4k^2-7k+2)\ge0$，第四段的差为
$2-B=4(2k-1)(k-1)\ge0$，并给出所述等号情形。
\end{proof}

\begin{remark}\label{cut:scope}
本定理闭合的是随阈值变化的一类非凸输入的真实三阶项，并给出统一四阶余项；
它不提供全体输入的三阶上界。$k=1$ 与未挖角正方形在三阶并列，
需要另算四阶项才能比较有限阈值优劣；本定理也未声称其四阶余项阶数不可提高。
\end{remark}
